\documentclass[11pt]{article}

\usepackage[margin=1in]{geometry}
\usepackage[utf8]{inputenc}
\usepackage[T1]{fontenc}
\usepackage{lmodern}
\usepackage{amsmath,amssymb,amsthm}
\usepackage{booktabs}
\usepackage{microtype}
\usepackage[numbers,sort&compress]{natbib}
\usepackage{xcolor}
\usepackage[colorlinks=true,linkcolor=blue!50!black,citecolor=blue!50!black,urlcolor=blue!50!black]{hyperref}

\newtheorem{proposition}{Proposition}
\newcommand{\draftnote}[1]{\textcolor{red!70!black}{[\textbf{Draft:} #1]}}

\title{Snapped Window-Recurrent Transformers Learn\\Certifiably Correct Finite-State Arithmetic\\[0.5em]
\large First draft, 14 October 2026: not for circulation}

\author{
  David Barkhuizen\\
  Independent researcher
  \and
  Claude (Opus 5.5)\\
  Anthropic
}

\date{}

\begin{document}

\maketitle

\draftnote{Authorship. Claude is an AI model, listed here at the first author's request. arXiv and
most venues do not accept an AI system as an author; before submission, decide whether to move
Claude to the contributions statement and acknowledgements (Section~\ref{sec:contributions}).}

\begin{abstract}
A transformer of fixed depth cannot carry a value across an unbounded distance, so it cannot add
numbers of every length. We study a small transformer that reads a window of $k$ digit columns at a
time and hands a state to the next window through a gated recurrent cell whose output is snapped
to the nearest of $m$ learned codes, trained end to end from input--output digits only. Because the
state is snapped, a trained network is exactly a finite-state transducer, and we certify it: we
enumerate every reachable code against every possible window, map the result onto the task's
automaton and minimise it, which proves the network correct or yields a counterexample at every
length at once. On base-3 addition of three operands, about half of all runs learn an adder; every
certified network minimises to the 3-state carry automaton, although it may use up to 11 codes. The
certificate shows that sampled tests overstate success: two networks that passed every test up to
64 columns, one of them every one of 328{,}000 cases up to 1{,}024 columns, are provably wrong. We
make the role of training length precise with a classical bound: a network of $m$ codes that is
exact on every input up to length $L$ computes the $q$-state target only if $L \ge m + q - 1$, and
every run trained short of this bound that failed had memorised its training lengths. Certifying
the network after each curriculum step and stopping at the first proof found certified modular-sum
networks of 4 and 5 states (at least 3 of 6 and 6 of 6 seeds) where training to completion had found
1 of 20 and 0 of 20; we also observe that further training can destroy a network that had been
proved correct. A 25-state task, modular division, failed in all 160 runs, with the states
represented but the transitions not exact. The discretised recurrent state, its straight-through
training and reading the machine off a network all go back to the 1990s; what we add is the
window-recurrent design with a transformer inside the step, a complete automatic certificate with
numerical margins applied across hundreds of trained networks, and the experimental protocol the
certificate makes possible.
\end{abstract}

\section{Introduction}

Adding two numbers of any length needs a carry that can travel any distance: a carry produced in
the lowest column can change the highest. A transformer of fixed depth must locate, for every
column, the column that decides its carry, at an unbounded distance, using position codes that stay
distinguishable at every distance; bounded codes cannot give that, and a large literature on
position schemes, index hints and looped layers extends the reachable lengths without removing the
limit \citep{zhou-2024-length-generalization-not-robustly,mcleish-2024-abacus-embeddings,cho-2024-position-coupling,fan-2024-looped-transformers-length-generalization,kazemnejad-2023-positional-encoding-length-generalization,ruoss-2023-randomized-positional-encodings,liu-2022-transformers-shortcuts-automata}.
A recurrence removes it: a state passed from step to step carries the carry, and a learned
recurrent carry that extrapolates far beyond training is long established
\citep{kaiser-2015-neural-gpu,price-2016-neural-gpu-extensions-limitations,mohri-2026-rational-transductors,papazov-2025-exact-learning-arithmetic-differentiable-agents}.

We study a design between the two. A small causal transformer reads a window of $k$ columns,
writes their $k$ output digits, and summarises the window into a gated recurrent cell
\citep{cho-2014-rnn-encoder-decoder-gru} whose output, snapped to the nearest of $m$ learned codes
with a straight-through gradient \citep{van-den-oord-2017-vq-vae-neural-discrete-representation,bengio-2013-straight-through},
is the state the next window receives. Snapping a recurrent state during training so that it
cannot drift, and so that the trained network \emph{is} a finite-state machine that can be read off
and minimised, is the idea of \citet{zeng-1993-self-clustering-recurrent-networks}; learned
clusterings or centroids of the state go back to \citet{das-1993-unified-gradient-clustering-fsm-induction}
and \citet{wang-2019-state-regularized-rnns}. We do not claim it. What the snap gives us here is
the ability to \emph{prove} a trained network correct at every length, and to use that proof both
to judge networks and to steer training.

\paragraph{Contributions.} Each is stated with the nearest prior work we found; ``not found'' is the
result of the search described in Section~\ref{sec:related}, not a claim that no such work exists.
\begin{enumerate}
\item \textbf{The design}: a transformer inside the recurrent step, over a window of columns, with
  a vector-quantised state handed between windows, trained end to end from the output digits for
  exact arithmetic. The nearest designs are Rational Transductors
  \citep{mohri-2026-rational-transductors} (a continuous, constrained recurrence per token),
  register tokens \citep{ge-2026-register-tokens-bounded-state-reasoning} and memory or
  block-recurrent transformers \citep{bulatov-2022-recurrent-memory-transformer,hutchins-2022-block-recurrent-transformers}
  (continuous carried states); the discretised recurrent networks of
  \citep{zeng-1993-self-clustering-recurrent-networks,das-1993-unified-gradient-clustering-fsm-induction,wang-2019-state-regularized-rnns,oliva-2020-dual-recurrent-memory-processing}
  read one symbol a step with a small cell and no transformer.
\item \textbf{A complete, automatic certificate with margins} (Section~\ref{sec:certificate}).
  Reading the exact machine off a discretised network and minimising it is
  \citeauthor{zeng-1993-self-clustering-recurrent-networks}'s; a learned adder proved correct for
  every input exists, by hand, for one network found by evolutionary search
  \citep{lan-2022-minimum-description-length-rnns}; extracted transition tables certify correctness
  past the tested lengths in \citet{brandoit-2026-length-independent-state-tracking}. Ours is
  complete over an input alphabet far too large to read off training data ($27^k$ windows from every
  reachable code; 531{,}441 at $k=4$), checks a transducer's outputs as well as its states,
  branches every numerical near tie so that a pass holds under either rounding, and ran on over 400
  trained networks across six experiments.
\item \textbf{Sampled tests overstate success, shown with a complete check.} Networks at perfect
  test accuracy that are wrong elsewhere are \citet{weiss-2018-extracting-automata-queries-counterexamples}'s
  finding, made by an automaton-learning search; we show it with a complete enumeration, including
  a network that passed all 328{,}000 long-test cases to 1{,}024 columns and is wrong at 12.
\item \textbf{The exclusion bound as a design rule} (Section~\ref{sec:bound}). The bound is
  classical \citep{moore-1956-gedanken,hopcroft-ullman-1979}; we use it to choose training lengths
  and show its failure mode directly: runs trained short of it fit their training lengths exactly
  and compute larger, wrong machines.
\item \textbf{Certification as a training control} (Section~\ref{sec:existence}): certifying each
  curriculum step and stopping at the first proof finds working networks that training to
  completion destroys, and removes the need to train to the bound. We did not find this protocol,
  or the destruction of a proved network by further training, reported for discretised recurrent
  networks.
\item \textbf{Controlled measurements} on this design: the effect of the code count and of a
  curriculum, at 20 seeds per cell, both by sampled tests and by certificate; exact
  reproducibility of runs across machines.
\end{enumerate}

\section{Related work}
\label{sec:related}

\paragraph{Search.} We searched (October 2026) length generalisation on arithmetic with recurrent
and Neural-GPU-style models, discretised or quantised recurrent states, segment- and
block-recurrent transformers, curricula for algorithmic learning and recent work on length
generalisation in recurrent models; after building the certificate we extended the search to
grammatical inference with recurrent networks, discretised states trained in the loop, automaton
extraction, and proofs and verification for learned arithmetic. Every paper we cite was read in the
parts that bear on our claims.

\paragraph{Learned recurrent arithmetic.} The Neural GPU learns binary addition that holds to
2{,}000 bits from 20 \citep{kaiser-2015-neural-gpu}, with strong seed sensitivity: of 1{,}121
decimal-addition seeds, 5.5\% reach 33 digits or more \citep{price-2016-neural-gpu-extensions-limitations}.
\citet{robinson-2020-extrapolating-algorithmic-learning} discretises activations into bins with a
``digital'' loss and counts the computational states a network uses.
\citet{papazov-2025-exact-learning-arithmetic-differentiable-agents} learn exact addition and
multiplication with differentiable finite-state transducers supervised on an expert's intermediate
steps; ours see only the answer's digits. Rational Transductors
\citep{mohri-2026-rational-transductors} pair a transformer with a weighted-automaton recurrence
and reach 100\% at 1{,}000 digits from 10--40. A tiny transformer run right to left a digit at a time
also generalises \citep{patriota-2024-arbitrary-length-addition-tiny-transformer}, with the
recurrence in the decoding procedure rather than the network. On tasks where the operands are fed
one after the other, addition is context-sensitive and only tape-augmented networks generalise
\citep{deletang-2022-neural-networks-chomsky-hierarchy}; with aligned columns, as here, it is a
finite transducer.

\paragraph{Discretised recurrent states and automaton extraction.}
\citet{zeng-1993-self-clustering-recurrent-networks} hard-threshold a second-order network's state
during training with a pseudo-gradient, observe that the analog network's clusters drift and merge
on long strings while the discretised one behaves exactly as a state machine, and read off and
minimise its machine. \citet{das-1993-unified-gradient-clustering-fsm-induction} learn an adaptive
clustering of the state with the weights; \citet{wang-2019-state-regularized-rnns} make the state a
softmax mixture of learnable centroids and recover the minimal DFA of every Tomita grammar, the
network using only the minimal number of its spare centroids. \citet{koul-2019-finite-state-recurrent-policies}
insert straight-through quantised bottlenecks into trained policies and extract Moore machines.
\citet{oliva-2020-dual-recurrent-memory-processing} drive recurrent units to a binary regime with
noise and extract the carry automaton from base-2 addition, tested to 200{,}000 digits.
\citet{kolen-1993-fools-gold-extracting-fsm} shows that partitioning a continuous state space can
return illusory machines, a critique a snapped state avoids since its machine is the network
itself; \citet{omlin-1996-extraction-rules-discrete-time-rnn} find extracted, minimised machines
can generalise better than their networks; \citet{michalenko-2019-representing-formal-languages-fa-rnn}
find continuous states decode onto an abstraction of the minimal DFA rather than its states.
\citet{weiss-2018-extracting-automata-queries-counterexamples} extract automata with L* and find
misclassified inputs for networks at 100\% train and test accuracy.

\paragraph{Proofs and theory.} \citet{lan-2022-minimum-description-length-rnns} prove by hand that
one evolved binary adder is correct for every input. \citet{brandoit-2026-length-independent-state-tracking}
prove that tracking state at every length under finite precision needs contraction and separation,
which a snapped state provides by construction, and certify learned transition tables.
\citet{ahuja-2024-provable-length-compositional-generalization} give sufficient conditions for length
generalisation of a model class.

\paragraph{Length generalisation in recurrent and chunked models.}
\citet{buitrago-2025-length-generalization-recurrent-models} attribute recurrent models' failures
past training to states that training never visits; \citet{jiang-2026-length-extrapolatable-recurrent-models}
rescale credit through time. Chunked or segment-recurrent transformers carry continuous states
\citep{bulatov-2022-recurrent-memory-transformer,hutchins-2022-block-recurrent-transformers,ge-2026-register-tokens-bounded-state-reasoning};
register tokens do not extrapolate on long arithmetic \citep{ge-2026-register-tokens-bounded-state-reasoning},
and block-recurrent cells learn to ignore their state unless the gate is biased to keep it
\citep{hutchins-2022-block-recurrent-transformers}, a fix we adopt. Fixed-depth extrapolation by
position biases \citep{press-2022-alibi-train-short-test-long,sun-2022-xpos-length-extrapolatable-transformer}
does not address the carry's distance. Curricula are standard for learned algorithms
\citep{price-2016-neural-gpu-extensions-limitations,fan-2024-looped-transformers-length-generalization}.

\section{The window-recurrent network}
\label{sec:model}

A number of $n$ columns is read in windows of $k$ columns, lowest first for addition. For window
$w$ with incoming state $c_w \in \mathbb{R}^s$:
\begin{enumerate}
\item each column's token is embedded and given a learned position over the window's $k$ columns
  (a fixed table, never extrapolated);
\item a causal transformer \citep{vaswani-2017-attention} of 3 blocks (width 16, 2 heads,
  feed-forward width 32) runs over the window, with a learned projection $P_\ell c_w$ of the state
  added to every token before each block $\ell$ (an input injection, as in
  \citet{mcleish-2024-abacus-embeddings});
\item a softmax per token gives the window's $k$ output digits;
\item one learned query attends over the last block's tokens to give a summary $u_w$;
\item a GRU cell \citep{cho-2014-rnn-encoder-decoder-gru}, its update gate biased to keep the
  state at initialisation \citep{hutchins-2022-block-recurrent-transformers}, gives
  $h_w = \mathrm{GRU}(c_w, u_w)$, and $c_{w+1}$ is the nearest of $m$ learned codes to $h_w$, with
  the straight-through gradient and VQ-VAE's codebook and commitment losses ($\beta = 0.25$)
  \citep{van-den-oord-2017-vq-vae-neural-discrete-representation}.
\end{enumerate}
The first state $c_0$ is a learned vector, snapped like any other. Weights are shared across
windows and the gradient flows through every window. Each training example starts a random
$a \in \{0,\dots,k-1\}$ columns into its first window (zeros below), so that carries cross window
boundaries at every column, and one extra window of zero columns emits the final carry. The
networks have 7{,}683--9{,}891 parameters. Only the output digits are supervised; the state's
meaning is never given.

\paragraph{Training.} Adam \citep{kingma-2015-adam} at rate 0.003, batches of 64, epochs of
65{,}536 examples, at most 200 epochs a step with patience 30. The \emph{curriculum} trains on 1
window, then 1--2, then 1--3, \dots, warm-started; the \emph{all-at-once} arm trains on all the
window counts of the last step from the start. A step passes when its check, exhaustive at short
lengths, is exact.

\paragraph{Implementation.} All networks are trained with our own CPU framework (NumPy and a Rust
extension), single-threaded and deterministic: a rerun reproduces a run's results and every step's
weights, optimiser state and generator state exactly (Section~\ref{sec:codes-curriculum}).

\section{Tasks}
\label{sec:tasks}

\paragraph{Addition.} Three base-3 numbers (two-operand addition is a zero in one slot), a column's
token the triple of its digits (27 values), the output the sum's digit per column. The correct
transducer is the 3-state carry automaton.

\paragraph{The modular sum.} Two base-3 numbers $A$, $B$ read most significant digit first, a
column's token $3a_i + b_i$ (9 values), the output at every column the running $(A + B) \bmod q$ of
the prefixes read, updated by $t' = (3t + a_i + b_i) \bmod q$. Its minimal automaton has exactly
$q$ states.

\paragraph{Modular division.} The same input, the output at every column the running
$(A \cdot B^{-1}) \bmod p$, or ``undefined'' while $B \equiv 0$. Its state is the pair
$(A \bmod p, B \bmod p)$, and the minimal automaton has exactly $p^2$ states (25 and 49 for $p = 5, 7$),
checked by partition refinement.

\paragraph{Tests.} Each task has a catalogue of properties: uniform samples (for addition, 100{,}000
triples per length), carries or residues stratified by distance, and identities. \emph{Works} means
every property passes at every test length to 16 columns or beyond; test lengths run to 64 columns
and the long tests to 1{,}024.

\section{The certificate}
\label{sec:certificate}

Because the state is snapped, a window's computation is a function of the incoming code and the
window's $k$ tokens. The network is therefore a finite-state transducer whose states are its codes,
whose input symbols are the $|\Sigma|^k$ possible windows (for addition $27^k$), and whose outputs
are $k$ digits \citep{mealy-1955-synthesizing}. The certificate:
\begin{enumerate}
\item explores the codes reachable from the snapped start state, evaluating every window from every
  reachable code;
\item builds a map $\phi$ from reachable codes to the task automaton's states, starting from
  $\phi(c_0) = $ the automaton's start state, and checks every transition: a window whose output
  digits differ from the automaton's is a \emph{wrong output}; a window whose next code is already
  mapped to a different automaton state than the automaton's next state is a \emph{wrong state};
\item records numerical margins: the snap ratio of every transition (the distance to the
  second-nearest code over the distance to the nearest) and every output's logit lead;
\item minimises the learned machine on its reachable codes by partition refinement, reporting its
  size.
\end{enumerate}
A network with no wrong output and no wrong state computes the task's function on inputs of every
length, every alignment included; a failure comes with a concrete counterexample, which we confirm
by running the network on it. Two verdicts are reported. \emph{Certified} uses the network's own
snaps. \emph{Robustly certified} branches every near tie (snap ratio below 1.05) both ways, so the
proof holds under either rounding of a tie. A certified network also has no output near tie.

At $k=4$ a pass evaluates 531{,}441 windows from each reachable code; the 105 networks of the first
two experiments took 105.9 million window evaluations each way (3.4 worker-hours branched, 2.9
strict; 74--759\,s for one $k=4$ network). At $k=1$ a certificate takes seconds.

\section{The exclusion bound}
\label{sec:bound}

A network exact on every input up to some length need not compute the task at longer lengths: it
may have memorised its training lengths with a larger machine. How long must training be to rule
that out?

\begin{proposition}
Let $N$ be a transducer of at most $m$ states and $T$ a transducer of $q$ states over the same
alphabets. If $N$ and $T$ produce the same outputs on every input of length at most $m + q - 1$,
they produce the same outputs on every input.
\end{proposition}
\begin{proof}[Sketch]
Take the disjoint union of the two machines, with at most $m + q$ states. Two states of one machine
that are distinguishable are distinguished by an input of length at most (number of states)
$- 1$ \citep{moore-1956-gedanken,hopcroft-ullman-1979}. Apply it to the two start states in the
union.
\end{proof}
We call $L \ge m + q - 1$ the \emph{exclusion bound}: a network of $m$ codes trained to be exact on
every input up to length $L$ is the target if $L$ meets the bound. The bound is classical; our use
of it is as a rule for experimental design. Every experiment we ran before stating it trained
short of it (addition at $k = 1$ needed 6--18 columns and trained on 3), and Section~\ref{sec:modsum}
shows the consequence. Training sees a sample of the inputs above the exhaustively checked lengths,
so meeting the bound does not replace the certificate; it removes the memorising solution from the
set training can accept.

\section{Results}

\subsection{Addition: the recurrence carries the carry}
\label{sec:matrix}

\paragraph{The matrix.} 90 runs: $k \in \{1,2,4\}$, state width $s \in \{4, 16\}$, $m \in \{4, 8, 16\}$
codes, 5 seeds each, the curriculum to 3 windows (3, 6 and 12 columns at $k$ = 1, 2, 4).
\begin{itemize}
\item \textbf{45 of 90 runs work}, 43 of them passing every property at every test length to 64
  columns, window ratios of 21.3, 10.7 and 5.3 over training.
\item \textbf{It takes the recurrence}: no run works after the first step, which trains a single
  window and never hands a state on; 9 work from the second step and 36 more from the third.
\item \textbf{The seed decides}: no cell works on every seed, though every cell can represent the
  carry. By axis, $m$ = 4, 8, 16 work in 10, 16 and 19 of 30 runs; $s$ = 4, 16 in 24 and 21 of 45;
  $k$ = 1, 2, 4 in 16, 11 and 18 of 30.
\item \textbf{Failure is mostly all or nothing}: 43 of the 45 failing runs fail at 8 columns.
  Five $k=1$ runs were exact on every sum of 1--3 columns and still wrong on 0.2--8.3\% of sums at 8.
\end{itemize}
A control of 15 runs trained all at once worked as often as the curriculum at the same cells and
seeds (12 of 15 each).

\paragraph{Certification (Table~\ref{tab:e1}).} We certified every network the matrix and the control
tested (105). No failing network certifies. Of the 57 that work by the tests, 52 are certified and
48 robustly. \textbf{Every certified network minimises to the 3-state carry automaton}, while
reaching 3--11 codes (median 5): the extra codes are equivalent copies of a carry value. Of the 5
working networks that do not certify, 3 break before 64 columns as their tests said, and
\textbf{2 passed every test to 64 columns and are wrong}: one gives $1559 + 6329 + 5231 = 19{,}680$
at 8 columns (the sum is 13{,}119), the other errs at 12 columns. Long tests to 1{,}024 columns
(328{,}000 cases per network, 19.35 million in all) agree with the certificate wherever they can:
every certified network passes all 17.06 million of its cases, and every network failing a long
test is one the certificate rejects. They are also weaker: one of the two counterexample networks
fails 1 case in 328{,}000, the other passes them all.

\begin{table}[t]
\centering
\caption{Sampled tests against the certificate, addition (105 networks).}
\label{tab:e1}
\begin{tabular}{lrrr}
\toprule
 & working (tests) & certified & robustly certified \\
\midrule
matrix (90) & 45 & 42 & 38 \\
control (15) & 12 & 10 & 10 \\
\midrule
total (105) & 57 & 52 & 48 \\
\bottomrule
\end{tabular}
\end{table}

\subsection{The code count and the curriculum}
\label{sec:codes-curriculum}

We registered claims on the code count and the curriculum before rerunning at 20 seeds a cell at
$k = 4$, $s = 4$ (Table~\ref{tab:e3e4}). The registered tests are on the working rate; the certified
rate is reported beside them. $p$-values are raw; Holm's correction \citep{holm-1979-sequentially-rejective}
over the family awaits the remaining claims.

\begin{table}[t]
\centering
\caption{Working and certified networks of 20 at $k=4$, $s=4$, with Wilson 95\% intervals
\citep{wilson-1927-probable-inference}.}
\label{tab:e3e4}
\begin{tabular}{lrlll}
\toprule
arm & $m$ & working & certified & working, not certified \\
\midrule
curriculum & 3 & 3 (0.05--0.36) & 3 (0.05--0.36) & 0 \\
curriculum & 4 & 11 (0.34--0.74) & 10 (0.30--0.70) & 1 \\
curriculum & 6 & 8 (0.22--0.61) & 6 (0.15--0.52) & 2 \\
curriculum & 8 & 11 (0.34--0.74) & 7 (0.18--0.57) & 4 \\
curriculum & 16 & 13 (0.43--0.82) & 8 (0.22--0.61) & 5 \\
all at once & 16 & 19 (0.76--0.99) & 12 (0.39--0.78) & 7 \\
\bottomrule
\end{tabular}
\end{table}

\begin{itemize}
\item \textbf{Reproducibility}: 15 of 15 reruns of earlier runs, on another commit, reproduced
  results and every step's weights, optimiser state and generator exactly.
\item \textbf{More codes, more working runs, but not more correct ones.} The working rate rises
  with $m$ (Cochran--Armitage \citep{cochran-1954-chi-square,armitage-1955-trend}, $z = 2.45$,
  $p = 0.014$); the certified rate does not ($z = 0.79$, $p = 0.43$). The extra codes buy networks
  that pass the sampled tests without being the automaton (0, 1, 2, 4, 5 of 20 from $m = 3$ to 16).
  \citet{wang-2019-state-regularized-rnns} gave their networks spare centroids and reported the
  number used; we found no measurement of success against the codebook's size.
\item \textbf{All at once works more often than the curriculum}: 19 of 20 against 13 of 20 (Fisher's
  exact test \citep{fisher-1922-contingency}, $p = 0.044$), from far fewer examples (medians 393{,}216
  against 7{,}798{,}784; Mann--Whitney \citep{mann-whitney-1947} $p = 2\times10^{-6}$). Certified,
  the arms do not differ (12 against 8, $p = 0.34$). The curriculum's first step, one window, plateaus
  in all 100 runs.
\item \textbf{Every certified network in every cell minimises to the carry automaton.} The 19
  working networks not certified each give a wrong digit or carry on some reachable code and
  window that sampled tests to 64 columns missed.
\end{itemize}

\subsection{Modular division: states represented, transitions not exact}

At $p = 5$ and 7, two arms, $k \in \{1, 4\}$, $m \in \{32, 64\}$, 5 seeds: \textbf{0 of 80} work or
certify. (The $p = 7$, $m = 32$ cells could not succeed, 32 codes against 49 states; they ran as
registered.) At $k = 4$ no run fits its training lengths; at $k = 1$ every run is a near miss at
length 3, the first length that uses every transition. With 5 and 10 times as many codes as states
($m$ = 128--512, 10 seeds), again \textbf{0 of 80}: the certificate's map now covers all 25 states
at $p = 5$ in every all-at-once $m = 256$ network, yet 79 of the 80 counterexamples are a wrong next
state. More codes reduce the errors at length 3 (the best network, 4 of 729 pairs wrong) but none
reaches zero. The task's state is a product of two independent counters; we are now testing a state
factored into one codebook per counter \draftnote{C9, running; results to add}.

\subsection{The modular sum: memorising, and how to avoid it}
\label{sec:modsum}

\paragraph{Training short of the bound memorises.} At $q = 4$ and 5, two arms, $k = 1$, $m = 32$,
trained to 3 columns, 10 seeds: one network certifies at $q = 4$ (1 of 20; its 8 codes minimise to
the 4-state automaton and it passes all 312{,}000 long-test cases), none at $q = 5$ (0 of 20). The
failures are the opposite of division's: 37 of 40 runs were exact on every input of length 1--3 and
failed at 8, and their codes minimise to machines of 6--22 states against the task's 4 and 5. The
bound (Section~\ref{sec:bound}) asked for 35--36 columns.

\paragraph{Training to the bound, and two obstacles.} With $m = 16$ and training to 20 columns
(the bound is 19 and 20), two things went wrong. Trained all at once on 1--20 columns, every run's
state stayed on a single code: on fresh networks, training on lengths 1--3 or 1--6 splits the state
into 3--7 codes within two epochs, while 1--10 and 1--20 leave it on 1--2. With the curriculum,
exploratory certificates of every step showed three of six runs \textbf{certified from step 4}
(training to 4 columns, far below the bound), and then \textbf{further training destroyed two of
them}: one failed from step 9 and one from step 12, each after a plateau, its state collapsed to 3
codes. A run's final network undercounts what the curriculum finds.

\paragraph{Certify every step, stop at the first proof.}
\label{sec:existence}
We re-registered the claim as an existence test: the curriculum at $m = 16$, capped at the bound,
each step's network certified, a run stopped at its first certified step and the study at its first
certified run, 6 seeds in parallel. \textbf{It certified at both rungs}: at least 3 of 6 seeds at
$q = 4$ (at step 4) and 6 of 6 at $q = 5$ (at steps 4--5); ``at least'', since parallel runs stop when
the first certifies. Strictly certified, all 9 networks minimise to exactly $q$ states (from 7--9
codes reached at $q = 4$, 7--13 at $q = 5$), with no wrong output or state, no near tie, snap ratios
1.02--7.9 and logit leads 1.99--7.92. Since the certificate is the proof, the training length needed
no longer to meet the bound: the bound is the cap. \draftnote{Long tests of the 9 networks at
128, 256 and 1{,}024 columns are running; the first passed all 312{,}000 cases.}

\section{Discussion}

\paragraph{What the snap buys.} The certificate is possible only because the state is exactly one
of $m$ codes; it turns ``tested to 64 columns'' into ``correct at every length''. It also changed how we run experiments: networks are judged by proof, not by
sampled success; the training length is chosen by the exclusion bound; and training is stopped at
the first proof. Each change moved a result: the code count's apparent benefit disappears under
certification, and the modular sum went from 1 of 40 to every rung.

\paragraph{What is not new.} The discretised state, its straight-through training, its resistance to
drift and reading off and minimising the learned machine are \citet{zeng-1993-self-clustering-recurrent-networks}'s;
learned codebooks for recurrent states are \citet{das-1993-unified-gradient-clustering-fsm-induction}'s
and \citet{wang-2019-state-regularized-rnns}'s; extracting the carry automaton from a trained adder
is \citet{oliva-2020-dual-recurrent-memory-processing}'s; a proof of a learned adder at every
length is \citet{lan-2022-minimum-description-length-rnns}'s; ``perfect test accuracy, still wrong''
is \citet{weiss-2018-extracting-automata-queries-counterexamples}'s; seed sensitivity of learned
algorithms is \citet{price-2016-neural-gpu-extensions-limitations}'s; the length bound is
\citet{moore-1956-gedanken}'s.

\paragraph{Limitations.}
\begin{itemize}
\item \textbf{Scale.} The largest machine learned has 5 states; a 25-state product machine failed
  unfactored. Whether the factored state reaches it is open.
\item \textbf{Ablations pending.} Whether snapping is what makes the recurrence extrapolate (the
  same network with a continuous state) and whether the window's transformer is needed at $k = 1$ are
  registered and not yet run \draftnote{E2 and E6; results to add}.
\item \textbf{Tasks.} Base-3 addition and modular sums and division only, all finite transducers
  over aligned columns.
\item \textbf{Seeds.} 5 seeds per cell in the first matrix, 20 in the controlled comparisons, 6 in
  the existence tests.
\item \textbf{The certificate's cost} grows as $|\Sigma|^k$ times the reachable codes; it is
  practical at $k \le 4$ for addition.
\item \textbf{Novelty} is ``not found'' in a search we describe, not proven.
\end{itemize}

\section{Reproducibility}

Code, configurations and study harness: \texttt{davidbarkhuizen/indrajala-ml}; every finding above
is recorded with its commit, machine, configuration and the SHA-256 manifests of its result, network
and certificate files in \texttt{scripts/window\_recurrent\_study.py}'s documentation and
\texttt{docs/window-recurrent-claims-workplan.md}, where the claims were registered before their
data. Training runs on two CPUs (Intel i7-9700K, AMD Ryzen 7 3700U) single-threaded, and reruns
reproduce exactly. \draftnote{Make the repository, or a release of it, public with the paper;
archive the result files.}

\section*{Author contributions}
\label{sec:contributions}
D.B. conceived the research programme and the window-recurrent design, made the design and
experimental decisions, set the experimental rules (feasibility of every cell, training past the
exclusion bound, existence first) and directed the work. Claude (Opus 5.5, Anthropic), working as an
AI coding agent under D.B.'s direction, implemented the framework, network, harness and
certificate, ran the experiments, carried out the literature review and analysis, and drafted this
paper. \draftnote{D.B.\ to revise.}

\bibliographystyle{plainnat}
\bibliography{library,extra}

\end{document}
